% ============================================================
% audit_table.tex
%
% A "system audit" table mapping the paper's claims to the evidence
% backing each one.  This is *the* canonical source of truth for which
% claims are PASS / PART / STUB / FAIL --- prose elsewhere in the paper
% should not contradict it.
%
% Status semantics (paste into the abstract or front-matter as well):
%   PASS  -- derivation complete or experiment confirms the claim to
%            stated precision.
%   PART  -- mechanism demonstrated, full quantitative match pending;
%            specific gaps documented in the evidence column.
%   STUB  -- analytical result with no numerical confirmation yet,
%            or planned experiment not yet run.
%   FAIL  -- tested and disconfirmed.
%
% Use \texttt{} for code/experiment names and file paths inside cells.
% Long cells are fine -- the column widths are tuned for prose.
% ============================================================

\label{tab:audit}

{\scriptsize
\setlength{\tabcolsep}{4pt}
\renewcommand{\arraystretch}{0.95}
\begin{longtable}{@{}p{2.6cm}p{3.8cm}p{3.2cm}p{4.4cm}p{1.0cm}@{}}
\caption{\small System audit: paper claims mapped to supporting
derivations and experiments.  Status semantics in the front-matter
(\texttt{PASS} / \texttt{PART} / \texttt{STUB} / \texttt{FAIL}).}
\label{tab:audit_caption} \\
\hline
\textbf{Claim} & \textbf{Mechanism} & \textbf{Continuum / formal limit} & \textbf{Evidence} & \textbf{Status} \\
\hline
\endfirsthead

\multicolumn{5}{l}{\textit{(Table~\ref{tab:audit}, continued from previous page)}} \\
\hline
\textbf{Claim} & \textbf{Mechanism} & \textbf{Continuum / formal limit} & \textbf{Evidence} & \textbf{Status} \\
\hline
\endhead

\hline
\multicolumn{5}{r}{\textit{(continued on next page)}} \\
\endfoot

\hline
\endlastfoot

% ── Kinematic channel ───────────────────────────────────────────────
Kinematic birefringence along $(1,1,-1)$
    & Closed-form uniaxial split from the spinor propagator: $M_\text{eff} = \tfrac43(\mathbb{I}-\hat{\mathbf{n}}_*\hat{\mathbf{n}}_*^T)$, so $|H_\text{RGB}(\mathbf{k})|^2 = 1-\tfrac43|\mathbf{k}_\perp|^2$. Group-velocity split $\Delta v = \tfrac23\sin\omega\,|\mathbf{k}|$ (angular law $\sin\alpha$, massive sector) and $\omega$-independent dispersion-curvature split $4/3$ vs $0$ (photon + massive)
    & Isotropic $c$ under $O_h$ averaging / $a\to 0$
    & \emph{Analytic:} eigenphase dispersion $\tan\theta = \tan(\omega/2)/|H_\text{RGB}|$ and group velocity in closed form; $|H_\text{RGB}|^2 = 1$ exactly along $(1,1,-1)$ at all orders. \texttt{exp\_02\_closed\_form\_split\_check}: structure-factor and dispersion identities to $<10^{-12}$, speed coefficients to $<10^{-5}$. \emph{Observability:} \texttt{exp\_01\_optical\_axis\_isotropy} ($129^3$, both engines) -- isotropic pulse flattens perpendicular to $(1,1,-1)$ (ratio $\to 0.01$), flat-axis alignment $=1.0000$, $\omega$-independent, survives exact A=1 renormalization.
    & \texttt{PASS} \\

% ── Gauge sector ────────────────────────────────────────────────────
Gauge-sector birefringence ($\mathbf{Q}$ eigenvalues $\{4,4,16\}$): uniaxial structure derived; observability \emph{open}
    & Induced photon action $\tfrac{ca^4}{2}\mathbf{F}^T\mathbf{Q}\mathbf{F}$; special direction Hodge-duals to the same axis $(1,1,-1)$. The bipartite RGB/CMY structure breaks $O_h$ to a uniaxial subgroup about $(1,1,-1)$ ($\mathbf{Q}$ fixed by $12/48$, RGB by $6/48$), so the action is genuinely uniaxially anisotropic
    & $O_h$-averaged Maxwell is the \emph{orientation-averaged} form, not what a fixed lattice's observables see; quantitative $\mathbf{Q}$ role is the coupling ratio $g_3^2/g_2^2 = 3/2$ (Paper~II)
    & \emph{Structure derived} (Paper~I App.~B; Paper~II). \emph{Observability OPEN:} no symmetry forces isotropy (the $O_h$-average is an external imposition the uniaxial dynamics do not justify); by the \texttt{exp\_01} analogy the anisotropy may survive in observables, where read literally it is a dimension-4, $O(1)$ effect excluded by $\sim$30--37 orders. The one possible escape -- electric/magnetic cancellation in the full dispersion (Paper~I did the magnetic sector only) -- is uncomputed. \texttt{exp\_03} (full E+B Peierls dynamics) is decisive. See \texttt{notes/gauge\_sector\_structural\_conclusion.md}.
    & \texttt{PART} \\

% ── Concordance ─────────────────────────────────────────────────────
Multi-channel concordance (P9)
    & Kinematic and gauge sectors inherit the \emph{same} optical axis $(1,1,-1)$ from the common bipartite geometry ($M_\text{eff}$ and $\mathbf{Q}$ special directions coincide)
    & N/A (cross-sector structural identity)
    & Axis-sharing is a \emph{structural} identity (derived): $M_\text{eff}$ and $\mathbf{Q}$ both inherit $(1,1,-1)$ from the same bipartite geometry. The kinematic dispersion is the one frontier-testable observational channel; the gauge sector's observable status is itself \emph{open} (see its row -- it may be unobservable, suppressed, or an excluded $O(1)$ effect, pending \texttt{exp\_03}). P9's strength as a multi-witness empirical claim is therefore hostage to that outcome.
    & \texttt{PART} \\

\hline
\end{longtable}
}
